\documentclass[11pt]{article}
\usepackage[margin=1in]{geometry}
\usepackage{amsmath,amssymb}
\usepackage[T1]{fontenc}
\usepackage{hyperref}
\usepackage{booktabs}
\usepackage{enumitem}
\hypersetup{colorlinks=true,linkcolor=blue,urlcolor=blue}
\setlist{nosep,leftmargin=*}
\emergencystretch=2em
\title{Mathematics of Northstar's Research and Portfolio Tools}
\author{Northstar project documentation}
\date{20 November 2026}
\begin{document}
\maketitle

\begin{abstract}
This paper describes the implemented calculations for technical research
decisions, trade guidance, portfolio risk, ledger-derived performance, and
historical and forward strategy evidence. It is intended for learning and
technical reproduction. The rules are deterministic summaries and review aids;
they do not forecast returns, establish suitability, or place trades. Each
quantity below names its units and sample convention because those choices
affect interpretation.
\end{abstract}

\section{Scope and notation}

Prices are daily closes unless stated otherwise. A return is a fraction in
equations unless a percent sign is shown; the interface often displays
percentage points. Let $P_t$ be a close, $V_t$ share volume, and $T_t=P_tV_t$
the observed traded value. Moving averages are arithmetic means over the
indicated number of observations. A session count is not a calendar-day count.
The app analyzes adjusted research closes; account valuation uses its separate
share-unit price series.

The code paths described here are implemented in
\path{decision_policy.py}, \path{indicators.py}, \path{trade_plan.py},
\path{portfolio_risk.py}, \path{account_performance.py}, and
\path{validation.py}. The accessible summary is \path{docs/MATHEMATICS.md};
detailed user-facing methodology and limitations remain in
\path{docs/REFERENCE.md}.

\section{Research score and decisions}

\subsection{Ranking score}

The overall score is a weighted average of five component scores $s_j$ on a
0--100 scale:
\[
 S=0.30s_{\rm trend}+0.25s_{\rm momentum}+0.20s_{\rm risk}
   +0.15s_{\rm timing}+0.10s_{\rm liquidity}.
\]
It is rounded to one decimal; ticker breaks ties. The trend score awards 25
points for each of $P_t>MA_{50,t}$, $P_t>MA_{200,t}$,
$MA_{50,t}>MA_{200,t}$, and positive three-month return. The momentum score is
\[
s_{\rm momentum}=\operatorname{clamp}_{[0,100]}\!\left(
50+25\frac{0.6R_3+0.4R_6}{\max(\sigma,10)}\right),
\]
where $R_3,R_6$ are percentage returns over 63 and 126 sessions and $\sigma$
is annualised volatility in percentage points. The risk component is
\[
s_{\rm risk}=\operatorname{clamp}_{[0,100]}(100-0.8\sigma+0.5D),
\]
where $D\leq0$ is the worst drawdown in the loaded history, in percent. Timing
uses RSI(14), MACD, and price extension. For RSI, let $g_t$ and $\ell_t$ be
the recursively smoothed mean positive and negative close changes: each starts
as the arithmetic mean of its first 14 changes and then updates as
$(13g_{t-1}+\max(\Delta P_t,0))/14$ and
$(13\ell_{t-1}+\max(-\Delta P_t,0))/14$. Then
\[
RSI_t=\begin{cases}50&g_t=\ell_t=0,\\100&\ell_t=0,\\
100-100/(1+g_t/\ell_t)&\text{otherwise.}\end{cases}
\]
EMA uses $E_t=\alpha x_t+(1-\alpha)E_{t-1}$,
$\alpha=2/(n+1)$, initialized at the first value. MACD is
$EMA_{12}-EMA_{26}$; its signal is the 9-period EMA of MACD. Timing starts at
90 points for $40\leq RSI\leq65$, 65 for $30\leq RSI<40$ or
$65<RSI<70$, and 35 otherwise; it adds 10 when MACD exceeds its signal (or
subtracts 10 otherwise), and subtracts 20 when the close is at least 15\% above
$MA_{50}$. The result is clamped to 0--100. The liquidity score is
\[
s_{\rm liquidity}=\operatorname{clamp}_{[0,100]}\!\left(
50+25\log_{10}\frac{\max(\bar T_{20},1)}{\max(L,1)}\right),
\]
where $\bar T_{20}$ is mean observed traded value and $L$ is the configured
NOK floor. Fewer than 15 valid turnover observations gives a zero score.
These components overlap and are not probabilities or independent evidence.

\subsection{Entry and held-position review}

Ranking order never overrides the entry gate. An entry-eligible stock must
have at least 200 fresh daily observations and matching currency, pass all
four trend conditions, have annualised volatility $<50\%$, close less than
15\% above $MA_{50}$, at least 15 valid traded-value observations in the last
20, and mean traded value at least $L$ (default NOK 250{,}000). A held stock
can qualify for an addition under the same screen. Portfolio constraints and
personal confirmation rules also apply before quantities are suggested.

For an owned holding, a technical sell review occurs when both
\[
P_t<MA_{200,t}\qquad\text{and}\qquad R_3<0.
\]
A saved lower review price triggers when the current quote is at or below it;
a saved upper review price triggers when the quote is at or above it. The
lower-level check takes precedence if both personal levels could apply.
Losing entry eligibility alone does not trigger a sale. The default review
quantity is 25\% of owned whole shares, rounded upward, subject to configured
portfolio reductions; zero disables that review quantity. Confirmation,
cooldown, cash, and allocation limits are guidance safeguards. No order is
transmitted.

\section{Trade-cost examples}

For budget $B$, price $P$, and one fixed buy fee $F_b$, whole shares are
\[
n=\max\left(0,\left\lfloor\frac{B-F_b}{P}\right\rfloor\right),
\qquad C_b=nP+F_b.
\]
The fee is charged only when $n>0$; if $n=0$, order cost and fee are both
zero. When $n>0$, fee drag is $100F_b/(nP)$ percent and fee-only break-even is
$P+F_b/n$. For $m$ shares sold at $P$, estimated proceeds are $mP-F_s$ only
if a sell fee $F_s$ is entered; otherwise the display remains gross at $mP$.
Taxes, spread, and execution slippage are not included in these personal
previews.

\paragraph{Checked example.} With $B=10{,}000$ NOK, $P=100$ NOK, and the
default $F_b=79$ NOK, $n=\lfloor9{,}921/100\rfloor=99$, order cost is
9{,}979 NOK, and 21 NOK remains. With average cost $C=100$ NOK/share, ten
shares, $F_b=79$ NOK and an entered $F_s=31$ NOK, the displayed break-even is
$100+(79+31)/10=111$ NOK/share. This is a fee estimate, not a recovery target.

\section{Portfolio risk}

For asset $i$, daily simple return over the exact interval $(a,b)$ is
\[
r_{i,(a,b)}=P_{i,b}/P_{i,a}-1.
\]
The portfolio view joins these interval keys exactly; it requires at least 60
common intervals for every holding and uses at most the latest 252. No missing
return is filled and an incomplete subset is not renormalized. USD assets are
converted using FX observations at both endpoints. Cash in NOK has zero return.

With current portfolio weights $w_i$ (including cash) and sample covariance
matrix $\widehat\Sigma$, the historical daily-rebalanced proxy has
\[
\widehat\sigma_p=\sqrt{252\,w^\top\widehat\Sigma w}.
\]
Variance contribution is $w_i\operatorname{Cov}(r_i,r_p)$; its share of total
variance divides by $\operatorname{Var}(r_p)$. Volatility contributions sum
to portfolio volatility; variance shares sum to 100\% when total variance is
positive. A hedge can contribute negatively. These are historical proxies
using today's weights, not account history.

Empirical 95\% VaR is the nonnegative magnitude of the interpolated 5th
percentile return. Expected shortfall averages exactly the worst 5\% of the
observed probability mass, including a fractional boundary observation.
Drawdown at time $t$ is $P_t/\max_{u\leq t}P_u-1$. Short samples and flat
series can leave statistics undefined.

Scenario shocks are arithmetic assumptions on current market value $V_i$.
For a stock shock $s$ and USD/NOK shock $f$, both in fractions,
\[
\Delta V_i=V_i\big((1+s_i)(1+f_i)-1\big),
\]
where $s_i=s$ for affected equities (or a selected classified sector), and
$f_i=f$ only for USD-denominated assets. The combined effect is multiplicative.
The default warning threshold is 20\%; it flags historical or scenario loss
estimates at or above the threshold. It is a review prompt, not a guaranteed
loss limit. Revenue exposure, hedges, costs, and taxes are not inferred.

\section{Ledger account performance}

Let $V_0,V_T$ be verified opening and ending account values and $F_k$ the
signed external cash flow on date $k$ (deposits positive, withdrawals
negative). Investment P\&L is
\[
\Pi=V_T-V_0-\sum_kF_k.
\]
Trades, fees, dividends, and interest are internal activity; recorded net
broker cash already reflects charges. Daily-linked return chains each
end-of-day flow-adjusted factor:
\[
G=\prod_t\frac{V_t-F_t}{V_{t-1}},\qquad R_{\rm linked}=G-1.
\]
This convention is an estimate because exports do not include exact flow
times. Modified Dietz is
\[
R_{\rm MD}=\frac{\Pi}{V_0+\sum_k w_kF_k},\qquad
w_k=\frac{\text{calendar days from flow }k\text{ to end}}
{\text{calendar days in period}}.
\]
If the weighted denominator is not positive, or flow timing makes linked
return undefined, the percentage is withheld while measurable NOK P\&L may
remain. Missing exact-date prices suppress the calculation. Explicit prior
close mode allows at most four calendar days and marks the result estimated.
The ledger must reconcile with current shares and tracked cash. Realized trade
P\&L uses verified weighted-average book cost and is not tax-lot accounting.

\section{Historical strategy testing}

The strategy experiment is a distinct model from personal trade plans. It
forms monthly decisions from bars through decision date $t$ and executes at
the next common session close. Holdings persist until an explicit technical
exit; qualifying entries fill vacant slots in score order, up to the chosen
cap. Equal-weight rebalance is an experimental overlay. Empty slots remain
cash. The baseline uses 20 slots and monthly decisions. Personal price levels,
confirmation, cooldowns, portfolio limits, and personalized share quantities
are not simulated.

For normalized equity $E_t$, a held portfolio with weights $w_{i,t}$ earns the
close-to-close gross return $\sum_iw_{i,t-1}(P_{i,t}/P_{i,t-1}-1)$ before
rebalancing costs. Each positive-value buy order incurs the stated 79 NOK
fixed fee; per-order experiment sale fees are configurable. Extra slippage is
charged in both directions in proportion to traded notional. Potential order
fees are reserved before sizing; unused reserves remain cash. Fixed per-order
fees are normalized by the configured initial capital. The benchmark is a same-stock equal
weight basket and, when available, the OSEBX gross-return index over exact
matching dates. Displayed excess is strategy net return minus index gross
return in percentage points, not risk-adjusted alpha. Strategy return uses
initial capital before the first execution's costs: a flat 100,000 NOK account
with one 79 NOK purchase fee returns $-0.079\%$, even against a flat index.
Index costs are not available, and neither side is terminally liquidated.

An investable comparator uses DNB OBX ETF (provider symbol OBXD.OL), which
tracks the 25-stock OBX basket rather than OSEBX GR. With capital $C$, entry
fee $F$, slippage fraction $s$, and adjusted entry and end prices $P_0,P_T$,
its ending wealth relative to initial capital is
\[
E_T^{fund}=(1-F/C)\frac{P_T}{P_0(1+s)}.
\]
Only dividend/split-adjusted NOK closes covering every matched session are
accepted. No dates are filled. Fund expenses are embedded in observed prices
and are not deducted twice. Fractional units, no terminal liquidation, taxes,
spreads, tracking differences and provider revisions limit comparability.
DNB's official fund page lists a 0.2\% management fee (checked 20 November
2026), not a verified historical fee schedule:
\url{https://www.dnb.no/en/markets/shares/exchange-traded-products/etf/dnb-obx}.

For a return sequence $r_t$, annualized volatility is
\[
\sigma_{ann}=\sqrt{252}\,s(r_t),
\]
where $s(r_t)$ is the sample standard deviation of fractional returns. Sharpe
at zero risk-free rate is $252\bar r/\sigma_{ann}$; Sortino is
$252\bar r/\sigma_{down}$, with
$\sigma_{down}=\sqrt{252\,\operatorname{mean}(\min(r_t,0)^2)}$. CAGR uses
elapsed years; drawdown is wealth relative to its running maximum. Reported
metrics require their stated sample counts.

The optional robustness grid compares 10/20 positions and 1/3 month
rebalance intervals. The chronological split is approximately 70\% training
and 30\% later holdout, with at least 20 sessions in each. Variant selection
uses training Sharpe at zero risk-free rate, then training return, with
deterministic ties. Holdings cross the boundary; holdout wealth is normalized
at the prior close. Extra 25/50 basis-point runs are sensitivity checks, not
independent confirmation. Repeated inspection or tuning against this
historical holdout contaminates any unseen-test claim.

Forward paper cohorts are a third, separate record. Each scan appends an
immutable snapshot; hypothetical entry uses the first later close, and returns
are shown only after a subsequent observation. Captured budget and fee
assumptions are retained. Cohorts overlap, are not a funded rebalanced
portfolio, and do not change primary advice. An experimental model affects
recommendations only after an explicit investor promotion.

\section{Limitations and reproducibility}

Historical experiments use today's configured universe and complete-history
assets. Delisted securities and reliable point-in-time fundamentals/news are
unavailable; adjusted prices can be revised. This creates survivorship,
selection, and data-revision bias. Complete-case filtering, fractional
historical allocations, simplified execution, omitted taxes, market impact,
and cash interest further limit inference. A matched index comparison does
not establish investable performance. Forward paper results are hypothetical
and have shorter, overlapping observation windows. None of these methods
establishes a forecast or investment result.

The maintained detailed explanations are in
\href{REFERENCE.md}{\texttt{docs/REFERENCE.md}}. Source implementations are
named in Section 1, and independent accounting examples and benchmark date
checks are maintained in \texttt{tests/test\_account\_reference.py}. This
paper describes the 20 November 2026 source snapshot; rerun its examples and
recompile after behavioral changes.

\end{document}
